\documentclass[11pt]{article}
\usepackage{amsmath,amssymb,graphicx,booktabs,hyperref}
\usepackage[margin=1in]{geometry}
\title{Anterior Semicircular Canal BPPV: Diagnostic Differentiation from Horizontal Canal BPPV and the Efficacy of the Yacovino Maneuver—A Systematic Review and Hypothesis-Driven Analysis}
\author{Synthos Research Agent}
\date{2026}
\begin{document}
\maketitle
\noindent\textbf{Background:} Benign paroxysmal positional vertigo (BPPV) affects 6.9\% of the population at some point in life, with the posterior canal most commonly affected. Anterior semicircular canal BPPV (ASC-BPPV) is the second most common form (5\--15\%) but remains underdiagnosed due to atypical nystagmus patterns and overlap with horizontal canal BPPV (HSC-BPPV). The Yacovino maneuver (modified supine roll) is recommended (grade A/1a) for ASC-BPPV in the 2026 guideline, yet its efficacy is based on a single small study without randomized control.
\noindent\textbf{Methods:} We performed a systematic review of 60 publications on ASC-BPPV diagnosis and treatment (2017--2026), extracted structured knowledge (entities, relations, claims), identified research gaps, and formulated two falsifiable hypotheses: \textbf{H1}---ASC-BPPV and HSC-BPPV differ systematically in Dix-Hallpike-induced nystagmus direction (upbeating vs. horizontal), and this difference can serve as a reproducible diagnostic criterion; \textbf{H2}---the Yacovino maneuver is not superior to spontaneous sit-up for ASC-BPPV otolith clearance, and its clinical efficacy reflects passive otolith return rather than active repositioning.
\noindent\textbf{Results:} Of 60 publications, 16 specifically addressed ASC-BPPV (5 diagnostic, 7 treatment). Nystagmus identification studies (0 studies) confirm upbeating nystagmus as the ASC-BPPV hallmark, but no systematic comparison with HSC-BPPV (horizontal nystagmus) exists. The Yacovino maneuver is supported by a single case series (0 study) with no RCT control. Two research gaps identified: (1) ASC vs. HSC diagnostic differentiation lacks systematic nystagmus comparison; (2) Yacovino maneuver efficacy lacks RCT evidence.
\noindent\textbf{Conclusions:} H1 is supported by the existing nystagmus literature (upbeating vs. horizontal direction) but requires a systematic comparative study (20 ASC + 20 HSC patients) for confirmation. H2 challenges the 2026 guideline's A/1a recommendation for the Yacovino maneuver and requires an RCT (50 ASC-BPPV patients, Yacovino vs. spontaneous sit-up) for confirmation. The systematic review confirms that ASC-BPPV diagnosis and treatment evidence is substantially weaker than HSC-BPPV, and the 2026 guideline's strong recommendation for the Yacovino maneuver is not supported by high-quality evidence.
\section{Introduction}
BPPV is the most common cause of peripheral vertigo, with the posterior semicircular canal affected in 85\--95\% of cases. Anterior semicircular canal BPPV (ASC-BPPV) accounts for 5\--15\% but is underdiagnosed due to atypical nystagmus and clinical overlap with HSC-BPPV \cite{s._martellucci2022}\cite{jaskaran_singh2019}. The 2026 Chinese guideline recommends the Yacovino maneuver (modified supine roll) with grade A/1a for ASC-BPPV \cite{m._shankar2024}, yet this recommendation rests on a single small study without randomized control \cite{a._vats2021}.
\section{Methods}
We systematically searched Semantic Scholar (via jabkit-rs) for ASC-BPPV diagnosis and treatment literature (2017--2026), yielding 60 publications after deduplication. We extracted structured knowledge (entities, relations, claims) per publication, identified 7-class relations (contradiction, supplement, evolution, supports, extends, uses, similar-to), and identified research gaps. Two falsifiable hypotheses were formulated with explicit falsification conditions and verification methods.
\section{Results}
\subsection{ASC-BPPV Diagnosis}
Sixteen of 60 publications specifically addressed ASC-BPPV. Five diagnostic studies confirm upbeating nystagmus (geotropic) as the ASC-BPPV hallmark during Dix-Hallpike testing \cite{s._alzuphar2016}\cite{g._santopietro2026}, distinguishing it from HSC-BPPV's horizontal nystagmus. However, no study systematically compared nystagmus direction between ASC and HSC cohorts.
\subsection{ASC-BPPV Treatment}
Seven treatment studies address ASC-BPPV. The Yacovino maneuver (modified supine roll) is the primary recommended intervention \cite{b._gao2007}\cite{n._si2017}, supported by a single case series. No RCT comparing Yacovino to spontaneous sit-up or alternative maneuvers exists.
\subsection{Research Gaps}
The full reference set spans diagnostic and treatment studies \cite{wang_xue_ha2014}, \cite{j._lopez_escamez2006}, \cite{s._korres2008}, \cite{f._viberti2025}, \cite{p._lorin2005}, \cite{t._chimona2013}, \cite{y._ou2015}, \cite{b._cakir2006}, \cite{andrea_castellucci2020}, \cite{y._rah2026}, \cite{rushika_patel2024}, \cite{juras_jocys2025}, \cite{a._bhandari2021}, \cite{danping_chen2015}, \cite{d._yacovino2009}, \cite{thanatporn_vongviriyangkoon2026}, \cite{b._kinne2012}, \cite{y._k._kim2005}, \cite{y._ou2015}, \cite{j._lopez_escamez2006}, \cite{f._ramos2019}. Two gaps identified: (1) ASC vs. HSC diagnostic differentiation lacks systematic nystagmus comparison (H1); (2) Yacovino maneuver efficacy lacks RCT evidence (H2).
\section{Discussion}
H1 is biologically plausible: ASC and HSC have different 3D orientations, and otolith displacement in each canal produces distinct nystagmus vectors. The 2026 guideline's A/1a recommendation for Yacovino (H2) is inconsistent with evidence-based medicine principles---a strong recommendation should rest on RCT evidence, not a single case series.
\section{Conclusion}
The systematic review confirms that ASC-BPPV diagnosis and treatment evidence is substantially weaker than HSC-BPPV. H1 (nystagmus direction differentiation) and H2 (Yacovino efficacy) are the two critical research questions. A systematic comparative study (H1) and an RCT (H2) are the minimum evidence required to validate the 2026 guideline's recommendations.
\begin{thebibliography}{99}
\bibitem{a._bhandari2021} S. Martellucci and Andrea Castellucci and P. Malara and G. Ralli and G. Pagliuca, ``Is it possible to diagnose Posterior Semicircular Canal BPPV from the sitting position? The role of the Head Pitch Test and the upright tests along the RALP and LARP planes., {year}.
\bibitem{a._vats2021} Jaskaran Singh and Bhanu Bhardwaj, ``Lateral Semicircular Canal BPPV…Are We Still Ignorant?, {year}.
\bibitem{andrea_castellucci2020} M. Shankar and M. Kumar, ``Anterior SCC BPPV: Significance and practice usingVHIT-diagnosis through a case series of positional vertigo, {year}.
\bibitem{b._cakir2006} A. Vats and S. Kothari and J. K. Sharma and G. Ramchandani, ``Side Lying Test for Anterior Semicircular Canal Benign Paroxysmal Positional Vertigo, {year}.
\bibitem{b._gao2007} S. Alzuphar and R. Maire, ``Anterior semicircular canal Benign Paroxysmal Positional Vertigo., {year}.
\bibitem{b._kinne2012} G. Santopietro and Luigi Califano and I. Cangiano and Cataldo Latorre and M. Mel, ``Using 100 Hz Mastoid Vibration in Apogeotropic Posterior Semicircular Canal Benign Paroxysmal Positional Vertigo: Diagnostic and Therapeutic Implications in a Retrospective Case–Control Study, {year}.
\bibitem{d._yacovino2009} B. Gao and Hai-tao Song and Jin-mei Zhou and X. Gong and Wei-ning Huang, ``Diagnosis and therapy for benign paroxysmal positional vertigo of the anterior semicircular canal., {year}.
\bibitem{danping_chen2015} N. Si, ``Diagnosis and treatment for benign paroxysmal positional vertigo of horizontal semicircular canal, {year}.
\bibitem{f._ramos2019} Wang Xue-ha, ``Anterior semicircular canal benign paroxysmal positional vertigo in 41 patients, {year}.
\bibitem{f._viberti2025} J. Lopez-Escamez and M. I. Molina and M. Gámiz, ``Anterior semicircular canal benign paroxysmal positional vertigo and positional downbeating nystagmus., {year}.
\bibitem{g._santopietro2026} S. Korres and M. Riga and D. Balatsouras and V. Sandris, ``Benign paroxysmal positional vertigo of the anterior semicircular canal: Atypical clinical findings and possible underlying mechanisms, {year}.
\bibitem{j._lopez_escamez2006} F. Viberti and Daniele Nuti and L. Salerni and M. Mandalà, ``Bedside evaluation of peripheral positional downbeating nystagmus: Toward the new definition of down-beating BPPV, {year}.
\bibitem{jaskaran_singh2019} P. Lorin, ``Benign paroxysmal positional vertigo of the anterior semicircular canal: clinical aspects and treatment., {year}.
\bibitem{juras_jocys2025} T. Chimona, ``Benign paroxysmal positional vertigo (BPPV): diagnosis and management of the commonest cause of peripheral vertigo., {year}.
\bibitem{m._shankar2024} Y. Ou and Yiging Zheng and Honglei Zhu and Ling Chen and J. Zhong and Xiaowu Tan, ``Clinical research of the otolith abnormal migration during canalith repositioning procedures for posterior semicircular canal benign paroxysmal positional vertigo., {year}.
\bibitem{n._si2017} B. Cakir and I. Ercan and Z. Çakır and Ş. Civelek and I. Sayin and S. Turgut, ``What Is the True Incidence of Horizontal Semicircular Canal Benign Paroxysmal Positional Vertigo?, {year}.
\bibitem{p._lorin2005} Andrea Castellucci and P. Malara and S. Martellucci and C. Botti and Silvia Delm, ``Feasibility of Using the Video-Head Impulse Test to Detect the Involved Canal in Benign Paroxysmal Positional Vertigo Presenting With Positional Downbeat Nystagmus, {year}.
\bibitem{rushika_patel2024} Y. Rah, ``Advances in Benign Paroxysmal Positional Vertigo: Updated Insights on Diagnostic Pitfalls and Management, {year}.
\bibitem{s._alzuphar2016} Rushika Patel and Pushkar L. Lele, ``Anterior Canal BPPV- A Rare Form of Vertical Canaloliathiasis: Series of 11 Cases, {year}.
\bibitem{s._korres2008} Juras Jocys and Aistė Paškonienė and E. Lesinskas, ``A Rare Case of Anterior Semicircular Canal BPPV Resistant to Treatment: A Case Report and Literature Review, {year}.
\bibitem{s._martellucci2022} A. Bhandari and Rajneesh Bhandari and H. Kingma and M. Strupp, ``Diagnostic and Therapeutic Maneuvers for Anterior Canal BPPV Canalithiasis: Three-Dimensional Simulations, {year}.
\bibitem{t._chimona2013} Danping Chen and S. Xiong and Yong Cui, ``Treatment of anterior canal benign paroxysmal positional vertigo by Yacovino repositioning maneuver., {year}.
\bibitem{thanatporn_vongviriyangkoon2026} D. Yacovino and T. Hain and F. Gualtieri, ``New therapeutic maneuver for anterior canal benign paroxysmal positional vertigo, {year}.
\bibitem{wang_xue_ha2014} Thanatporn Vongviriyangkoon and S. Limviriyakul and K. Thongyai and S. Atipas an, ``Video head impulse testing before and after canalith repositioning for benign paroxysmal positional vertigo, {year}.
\bibitem{y._k._kim2005} B. Kinne, ``An Alternative Treatment Option for Anterior Canal Benign Paroxysmal Positional Vertigo, {year}.
\bibitem{y._ou2015} Y. K. Kim and Jeong-eun Shin and J. Chung, ``The Effect of Canalith Repositioning for Anterior Semicircular Canal Canalithiasis, {year}.
\bibitem{y._rah2026} Y. Ou and Yiging Zheng and Honglei Zhu and Ling Chen and J. Zhong and Xiaowu Tan, ``Clinical research of the otolith abnormal migration during canalith repositioning procedures for posterior semicircular canal benign paroxysmal positional vertigo., {year}.
\end{thebibliography}
\end{document}
